\section{Implementation}
\label{sec:implementation}

\subsection{SGLang integration}

\method{} is implemented in a SGLang fork as a controller-owned KV pool
beside the ordinary request cache. The reuse arm publishes a dynamic
manifest of sources and cases; the Dense arm uses the same generate path
with copy disabled. Ordinary prefix reuse is compiled off in both arms
so the contrast is island copy versus Dense prefill, not island copy
versus radix.

Each source request writes selected $K/V$ into controller-owned slots.
Materialization applies \texttt{pre\_rotate\_delta} to $K$ only, then
records the applied delta in a per-source map. At the target the leftover
rotate is
$\Delta_{\mathrm{res}}=(t_{\mathrm{start}}-s_{\mathrm{start}})-\Delta_{\mathrm{source}}$.
The admitted campaign sets $\Delta_{\mathrm{source}}$ to the logical
shift, so $\Delta_{\mathrm{res}}=0$ and the copy path does not rotate
again. Before a copy, the controller checks source token-ID hashes and
full target coverage. Hash mismatch, an unowned gap, allocation failure,
or a copy that is not mechanically valid raises a fallback event and the
span is Dense-recomputed. The server ledger records copy events, fallback events, and the
applied pre-rotate delta. Mechanical success is read from that ledger,
not inferred from latency. Copy tiles stay at the engine default of
256 tokens.

Listing~\ref{list:policy-config} is the conceptual manifest shape.
Ordinary prefix reuse (M1) and prefetch (M3) are off in job 137185.
A prefetch miss copies from the still-owned handle (M2) instead of
Dense-recomputing the island.

\begin{lstlisting}[language=Python, caption={Reuse manifest (conceptual)}, label={list:policy-config}, breaklines=false, columns=fullflexible, basicstyle=\ttfamily\footnotesize]
manifest = {
  "prefix_reuse": False,  # M1 off
  "prefetch": False,      # M3 off
  "sources": [{
    "id": "read:mod.py@v3", "start": 1842,
    "length": 1109, "pre_rotate_delta": -1250,
  }],
  "cases": [{
    "target_start": 592, "length": 1109,
    "source_id": "read:mod.py@v3",
  }],
}
\end{lstlisting}

\subsection{Exact-prompt replay}

The speed campaign does not re-sample the live agent. It reconstructs
the exact target token IDs from the frozen trajectories with the same
chat template and rolling-6 compaction, then replays those IDs. Decode
is one token. Each group has one warmup and three measured rounds. The
reuse arm materializes sources once per group; the Dense arm does not.
Cache-ready TTFT is the target generate time after sources already
exist, matching the protocol used by CacheBlend and KVCOMM.

Server processes restart every 20 groups to bound long-run memory drift.
A group that fails mid-flight is retried after a fresh server. Those
guards are operational; they are not part of the reuse policy.

\subsection{Model and hardware}

The headline model is Qwen2.5-Coder-7B-Instruct. The campaign ran as
Slurm job 137185 on a compute GPU (not a debug or login node).
Temperature is 0. The same protocol on Qwen3-Coder 30B-A3B Instruct
(AWQ 4-bit, job 96092) is Appendix~\ref{app:30b}; we do not mix that
table with the 7B headline or with another model family's
official-Accuracy table.

A PLAN group is hashes, $(s,t,L)$, and $\Delta\neq 0$. Length, $N$-use,
island-count, and repo slices are derived from the frozen
\texttt{dense.json}/\texttt{reuse.json} pair; they are not a second GPU
campaign. Porting to another tokenizer means decoding the island text and
re-finding the span. Replaying 30B token ids on a 7B vocabulary is invalid.
